\section{Introduction} \label{sec:intro}

Autonomous robotic systems are a prerequisite for sustainable planetary exploration, particularly where communication latency and environmental hazards preclude continuous human supervision. The forthcoming establishment of permanent lunar infrastructures and the exploitation of in-situ resources require robotic platforms capable of performing complex manipulation and mobility tasks in environments that are not only remote but also extremely unforgiving. Missions such as NASA’s Artemis~\citep{nasa_artemis_2025} and China’s Chang’e~\citep{JIA2018207} programs highlight this growing reliance on robotic systems for long-duration lunar surface operations, emphasizing autonomy as a cornerstone for sustained human and robotic presence on extraterrestrial bodies.

To date, planetary exploration has relied primarily on wheeled rovers, which have demonstrated remarkable robustness and operational simplicity across decades of missions, from Sojourner and Spirit to Curiosity and Perseverance~\citep{mars-rovers}. However, despite advances in wheel design and suspension systems~\citep{Johnson2015, Zheng2018}, their performance degrades drastically on steep slopes, uneven surfaces, and soft regolith: conditions that dominate the lunar and martian landscapes. The Spirit rover’s immobilization in soft sand is a well-known example of such limitations~\citep{spirit}. These constraints have spurred growing interest in legged mobility for planetary exploration, where adaptability, stability, and terrain compliance are critical to mission success.

Legged robots, inspired by biological systems, offer superior versatility to adapt their gait and posture to the terrain morphology~\citep{Kolvenbach2022}. In contrast to wheeled systems, legged platforms can maintain contact stability over discontinuous surfaces \citep{Chong2024RiskyTerrain}, climb obstacles exceeding their body height \citep{vogel2025robustladderclimbingquadrupedal}, and traverse slopes beyond 15--20 degrees~\citep{Valsecchi2023}. Early research prototypes such as SpaceClimber~\citep{5756948}, ATHLETE~\citep{Wilcox2007}, and SpaceBok~\citep{8794136} have demonstrated the feasibility of multi-legged and dynamic quadrupedal mobility in reduced-gravity and granular conditions. More recent studies on steep terrain locomotion, such as the crawling gait adaptation of ANYmal~\citep{Valsecchi2023}, have shown that leg reconfiguration and new stability metrics can significantly enhance climbing robustness and reduce tumbling risk on sandy slopes up to 25° inclination. Together, these developments indicate a paradigm shift from passive, pre-programmed exploration vehicles toward adaptive, embodied agents capable of intelligent locomotion across unstructured extraterrestrial terrain \citep{Kolvenbach2024}.

Beyond locomotion, upcoming lunar operations will demand more than passive observation; they require interaction with the environment through sampling, assembly, and maintenance \citep{Walther2024InSituBoulders}. Quadrupedal mobile manipulators, which combine the mobility of legged locomotion with the dexterity of articulated arms, represent a promising solution to these dual requirements (see \cref{fig:intro}). However, integrating locomotion and manipulation into a unified control framework introduces a high-dimensional, nonlinear control problem where balance, force distribution, and task execution must be coordinated in real time. Achieving this level of autonomy requires methods that can learn complex dynamic behaviours while respecting explicit safety and stability limits.

\begin{figure}
    \centering
    \includegraphics[width=\columnwidth]{figures/intro3.png}
    \caption{Legged Mobile Manipulator in Lunar environment used in our work. The proposed system combines quadrupedal locomotion with a 6-DoF manipulator for whole-body control under low-gravity and irregular terrain conditions.}
    \label{fig:intro}
\end{figure}

Reinforcement Learning (RL) \citep{sutton2018reinforcement} has emerged as a powerful paradigm for enabling robust, adaptive control in high-dimensional robotic systems. RL-based controllers have achieved state-of-the-art performance in quadrupedal and humanoid locomotion~\citep{HoellerAnymalParkour, MikiWildAnymal, RLHumanoid}, dexterous manipulation~\citep{openai2019solvingrubikscuberobot}, and integrated navigation~\citep{rudin2022advancedskillslearninglocomotion, shah2022viking}. In particular, deep RL has been shown to produce agile, perception-driven policies that generalise across varied terrains and physical parameters~\citep{MikiWildAnymal}.

The progress made by RL-based legged platforms has naturally motivated their extension to loco-manipulation tasks. Nevertheless, most existing works simplify the problem by decoupling locomotion and manipulation. This separation is handled either explicitly via distinct controllers or modes dedicated to each capability \citep{Ma2022Locomanip}, or implicitly, through reward formulations that treat both objectives independently, allowing the policy to optimise them without promoting true whole-body coordination \citep{fu2022deep}. This separation limits performance, as it prevents the system from exploiting whole-body momentum and reachability to execute aggressive or complex tasks. More recent efforts have begun to unify these modalities and address whole-body control directly, enabling the robot’s legs and arms to be jointly optimised for coordinated motion and interaction. For example, \citep{portela2025wholebodyendeffectorposetracking} learns a 6D end-effector tracking controller for a legged manipulator, but only under static stance conditions, delegating locomotion to a separate learned module. Other approaches achieve simultaneous locomotion and manipulation \citep{ha2024umilegs, He2024}, yet rely on an additional dedicated manipulation policy to track arm trajectories, increasing system complexity and reducing end-to-end generality.

While these methods demonstrate impressive versatility in terrestrial environments, they generally lack formal mechanisms to ensure safety, such as explicit limits on body orientation, contact forces, actuation effort, or power consumption. Such omissions are acceptable in laboratory settings, where failures are recoverable, but they are prohibitive in planetary contexts where failures are catastrophic and recovery is impossible.

Operating on the lunar surface further compounds these challenges. The Moon’s reduced gravity (approximately one-sixth of Earth’s), its irregular and discontinuous topography, and the presence of soft, cohesion-less regolith create a domain where even small control instabilities can escalate into mission-threatening events \citep{Kolvenbach2022}. Moreover, communication latency, severely limited bandwidth, reduced computational capabilities, and stringent power budgets preclude continuous teleoperation or high-frequency supervisory intervention. As a result, autonomous systems must exhibit inherently self-preserving behaviour maintaining stability, safety, and efficiency despite uncertainties in terrain morphology, contact dynamics, and environmental disturbances, all this while being deployed within limited computational resources. These constraints demand control frameworks that embed safety, robustness, and adaptability directly into the learning process, rather than layering them on post hoc.

Constrained Reinforcement Learning (CRL) provides a principled avenue for addressing these challenges. Unlike standard RL, which optimises a reward function alone, CRL incorporates explicit safety and performance constraints that guide policy learning to avoid violating them~\citep{achiam2017constrained}. This yields policies that are optimized subject to operational limits expressed in physical units, rather than policies in which safety is traded off implicitly through reward weights. CRL has recently demonstrated promising results in terrestrial legged locomotion~\citep{LeeCRL, KimNotOnlyRewards, chane2024cat, chanesane2024solo} and emerging loco-manipulation tasks~\citep{dadiotis2025dynamicobjectgoalpushing, MaBadminton}. Its application to planetary robotics remains limited, however, and in particular no prior work outside our own preliminary study \citep{belmonte2025isparo} has combined constrained whole-body loco-manipulation with the gravitational and substrate conditions of the lunar surface; \cref{sec:related} surveys this literature in detail. The integration of CRL into lunar robotic systems therefore represents a critical step toward achieving autonomy that is both high-performing and inherently safe.

This work introduces a CRL framework for autonomous legged mobile manipulators designed specifically for lunar surface operations. Our framework integrates locomotion and manipulation into a unified whole-body control policy while embedding safety-critical constraints such as collision avoidance, dynamic stability, and power efficiency directly within the learning process. By combining RL-based behaviour with explicit constraint enforcement, the approach seeks a favourable trade-off between task performance and adherence to the operational envelope. We train and evaluate our learning framework in simulation by recreating challenging scenarios of lunar surface operations, such as rough, changing terrain topography, as well as lunar soil effects during contact phases of the locomotion gait.

A preliminary version of this framework was presented in ~\citep{belmonte2025isparo}, where a single policy was evaluated on one rough-terrain morphology under reduced gravity. The present article revises the reward formulation, most notably through the introduction of an explicit stability term and a generalised treatment of reward regularization under termination-based constraints; reformulates the constraint set around terrain interaction rather than generic postural limits; develops the environment model into a lunar-analog suite spanning multiple terrain morphologies, difficulty levels and regolith densities; and reports an entirely new experimental campaign, including a constraint ablation over separately trained policies. The resulting tracking errors improve on \citep{belmonte2025isparo} while being obtained under considerably more demanding terrain and contact conditions.

Our contributions are as follows.

\begin{itemize}
    \item \textbf{A whole-body loco-manipulation policy for task-space 6D pose tracking under lunar conditions.} We formulate end-effector pose tracking with the target specified in an inertial task frame, which makes coordinated base and arm motion necessary rather than optional, and solve it with a single end-to-end policy without a separate manipulation module or mode switch. We introduce a reward formulation in which the transition between locomotion-dominated and manipulation-dominated behaviour arises from the geometry of a saturating base-proximity term rather than from a gating function or state machine, and in which error-adaptive kernel widths sustain a usable tracking gradient across the three orders of magnitude of position error spanned by an episode.

    \item \textbf{A stability reward grounded in reduced-gravity dynamics.} We incorporate the Gravito-Inertial Inclination Margin (GIIM) as a dense training signal rather than an evaluation metric, using a smooth mean-over-faces surrogate in place of the non-differentiable minimum. We show that under lunar gravity the inertial contribution to the net acceleration acting on the body is a substantial fraction of the gravitational term, so that a gravito-inertial stability criterion is a requirement of the operating regime rather than a refinement of static support-polygon criteria.

    \item \textbf{A design rule for reward regularization under termination-based constraints.} We show that the multiplicative modulation introduced by our constrained learning algorithm of choice systematically undermines penalty-form regularizers whose penalised quantities are dynamically coupled to the constrained ones, attenuating the penalty precisely when a violation occurs and vanishing entirely in the hard-constraint limit. We establish that expressing every regularizer as a bounded non-negative reward preserves the policy ordering it induces under the modulation, and apply this rule consistently across the power, action smoothness and stability terms.

    \item \textbf{A lunar-analog environment model and a systematic robustness characterisation.} We construct a simulation environment spanning three terrain morphologies at five difficulty levels under lunar gravity, coupled with a model of regolith contact mechanics calibrated to the JSC-1a simulant and spanning loose, nominal and dense densities. Contact compliance is withheld from training, so that evaluation additionally measures robustness to this class of model mismatch. We characterise tracking performance across the resulting terrain-difficulty-soil combinations and quantify the contribution of individual constraints through an ablation over separately trained policies.
\end{itemize}

Together, these contributions address the conjunction of unified whole-body control, explicit safety enforcement, and planetary environmental conditions that \cref{tab:related} shows to be unaddressed in the existing literature, and establish a quantitative baseline for safe autonomous loco-manipulation on the lunar surface.